\section{Data, events and targets}\label{sec:data}

\paragraph{Collectors and archives.}
We use the same ten public route collectors in every year.
Seven are RIPE RIS collectors~\cite{ripe_ris}: rrc00 and rrc25 (multihop), rrc01 (London), rrc03 (Amsterdam), rrc12 (Frankfurt), rrc15 (S\~ao Paulo) and rrc20 (Zurich).
Three are RouteViews collectors~\cite{routeviews}: route-views2, route-views3 and route-views.linx.
We chose them with a greedy heuristic over the BGPKIT peer statistics of 2024-06-15~\cite{bgpkit}.
Each step adds the collector with the most full-feed peer ASes not yet covered, 235 in total.

We stream all available update archives from January 2022 to August 2026, 2.8--3.9\,TB per year.
The R2 forecasts use the same collectors from 1 September to 16 October 2026, and the beacons run to September 2026.
% src (round 7, criterion M17): R2 window tab:month-boundary header (1 Sep--16 Oct 2026); beacon period sections/07-beacons.tex and
%   "all 57 months, January 2022 to September 2026" (sections/07-beacons.tex, Trend 2022--2026).
We keep per-prefix origin changes, new AS adjacencies and path triples.
% round 6 (panel-r6 S9 / T6, panel-r5 NEW-14): "first-seen" now names only the leaker class of sections/08-forecast.tex.
% src: docs/paper/corrections-2026-10-16/vp-selection-verification.json; ingestion-manifest-verification.json
In 2024, a wider RIS pass also found 317 selected origin changes at 11 to 24 collectors, but every replay uses ten.
% layout (round 5, allowlist revision): this sentence moved here from "Event identification", which now holds the allowlist thresholds (six sentences at most).
% src (round 5, review-r5 W4.2 / item 4, panel-r5 W6(b) / X9): docs/paper/corrections-2026-10-16/discovery-collector-scope.json
%   years.2024: stratified 18 + labelled 248 + rpki_invalid_sample 51 = 317 discovered_at_more_than_ten_collectors, max_collectors 24;
%   every other year 0 and max 10. docs/paper/corrections-2026-10-16/README.md: "The 2024 discovery inventory also includes the wider RIS
%   beacon-week pass ... All replay footprints use the declared ten collectors." The origin recall comes from the ten pass collectors
%   (docs/results/2026-09-04-inventory-2024.md, identifier v2 "on the per-prefix tables of the 10 pass collectors").

\paragraph{Event identification.}
We identify events without human labels.
An \emph{origin change} starts when the observed origin AS of a prefix changes at two or more collectors within a five-minute bin.
A retrospective allowlist removes pairs of prefix and new origin that change in many five-minute bins of their calendar month.
% src (round 4, panel-r4 Q19 / review-r4 W5 "allowlist window"): src/bgbench/identifier2.py dedupe_origin_events (pairs with >= f bins in
%   the analysis window are whitelisted) and run_events (window = start..end of the call); jobs/inventory_year.sh (one call per calendar month,
%   --flapper-bins 50); docs/results/2026-09-04-inventory-2024.md ("same-window flapper whitelist F = 50 bins ... month by month"; 7,827 labelled
%   hijacks whitelisted away in 2024); its monthly recall table equals docs/paper/corrections-2026-10-16/identifier-summary-verification.json
%   (e.g. January 10--49 peers 0.690, 50+ 0.81 of 177), the source of the printed 2024 recall. docs/paper/corrections-2026-10-16/README.md:
%   "The origin-change whitelist is computed over the analysis window, rather than being an online rolling detector."
%   The threshold of the 2022, 2023, 2025 and 2026 runs, "not established" in round 4, is now established (round 5, next src comment).
A revision analysis finds a threshold of 50 bins in 2024, but 200--221 per month, or 50 per week, in the other years.
A flat 50 would remove 22.5, 18.4, 9.8 and 14.3\,\% of the selected origin changes of 2022, 2023, 2025 and 2026.
% src (round 5, panel-r5 NEW-9 / X10, A-R1(a), domain B2; review-r5 W4, item 4): paper/generated/allowlist_by_year.json
%   (scripts/paper_allowlist_by_year.py, first commit b644cf6b 2026-10-31 17:47 UTC): by_year.2024 {rule flat, thresholds_bins [50]},
%   by_year.{2022,2023,2025,2026} {rule per_week_scaled, thresholds_bins [200, 214, 221]}, confidence "established" in every year; method
%   "max over kept pairs of sum(n_bins) < T <= min(bins) of flapper_whitelist.parquet"; years.Y.monthly[*].threshold_bins_applied; run
%   provenance years.Y.run_notes / run_commits_contain_86b9b4b2 (2024 written 16--24 min before 86b9b4b2, the other years after it).
%   Flat 50: years.Y.downstream.replay_selected_in_pairs_bins_ge_F_lt_scaled / replay_selection_rows (= eventset_rows = release_eventset_rows):
%   2022 1,120 / 4,977 = 22.5 %, 2023 2,044 / 11,107 = 18.4 %, 2025 2,496 / 25,579 = 9.8 %, 2026 2,465 / 17,248 = 14.3 %; 2024 0 / 20,959.
%   The replay selection = catalogue population + label-free stratum (tab:selection: 2022 1,980 + 2,997 = 4,977).
%   Not printed: the same-run recall at 10--49 peers of the scaled years (validation_threshold_note: 2023 0.78--0.86 over 9 months,
%   2026 0.81--0.88 over 8 months, none for 2022 and 2025), "not a controlled comparison: label coverage and year differ".
Monthly recall of the stricter 2024 identifier is 0.66--0.73 against Cloudflare Radar's hijack catalogue~\cite{cloudflare2023hijack,cloudflare_radar_bgp} for events that ten or more of its peers saw.
% src (round 5): allowlist_by_year.json paper_2024_recall_provenance (recall_10_49_range_2024 [0.662, 0.727]; 12 of 12 pinned evaluation files equal
%   the flat F = 50 outputs) and validation_threshold_note.same_threshold_as_other_years false: the 2024 recall scores the stricter identifier.

A \emph{route leak} candidate is a path triple $(x, a, b)$ that is new in the current month.
It must violate valley-freeness under CAIDA's AS relationships~\cite{luckie2013relationships,caida_asrel}, where AS $a$ learns the route from $b$ and exports it to $x$.
New triples that violate valley-freeness match 8--11\,\% of the Radar leak labels per month in 2024.
All new triples match 17--22\,\%, in checks at 26 collectors.
% src (round 5): docs/paper/corrections-2026-10-16/identifier-summary-verification.json months[*].leak.collectors = 26 in all 12 months of 2024.
The identifier thus misses leaks on reused triples.

We do not analyse withdrawal waves, because only 56 of 301 have the recent reference RIB at all ten collectors that their path loss needs.
% src: docs/paper/corrections-2026-10-16/identifier-summary-verification.json; withdrawal/withdrawal-baseline-verification.json (50 waves 2024 + 6 in 2025 supported of 301)

\paragraph{Replay selections.}
The \emph{catalogue population} joins two parts (\Cref{tab:selection}).
The first holds the identified origin changes that match a Radar label visible at ten or more peers.
The second is a stratified sample of all identified origin changes.
Radar labels begin in 2023, so the 2022 set holds the stratified sample only.
A separate \emph{label-free invalid stratum} holds RPKI-invalid origin changes seen at ten or more peers and three or more collectors, selected without any catalogue label.
It tests how much the catalogue shapes the invalid class.

The leak selection holds Radar matches seen at three or more collectors plus the largest candidates by message volume.
Without Radar leak labels, 2022 and 2023 hold the top 396 candidates by collectors and messages.
\emph{Usable} events have complete archive coverage and two hours of follow-up.

\begin{table}[t]
\footnotesize
\centering
\caption{Replay selections by year: selected candidates and usable events. Cat.: catalogue population (2022: stratified sample only). Str.: label-free invalid stratum. 2026: January to August.}
\label{tab:selection}
\setlength{\tabcolsep}{3pt}
\begin{tabular}{@{}lrrrrrr@{}}
\toprule
 & \multicolumn{4}{c}{Origin changes} & \multicolumn{2}{c}{Route leaks} \\
\cmidrule(lr){2-5}\cmidrule(l){6-7}
 & cat. & usable & str. & usable & selected & usable \\
\midrule
2022 & 1,980 & 1,963 & 2,997 & 2,975 & 396 & 396 \\
2023 & 8,240 & 8,200 & 2,867 & 2,865 & 396 & 396 \\
2024 & 18,367 & 18,350 & 2,592 & 2,590 & 1,602 & 1,602 \\
2025 & 22,933 & 22,921 & 2,646 & 2,630 & 8,437 & 8,412 \\
2026 & 14,506 & 14,500 & 2,742 & 2,740 & 6,416 & 6,411 \\
\midrule
total & 66,026 & 65,934 & 13,844 & 13,800 & 17,247 & 17,217 \\
\bottomrule
\end{tabular}
\end{table}
% src: paper/generated/propagation-facts.json, summed by script: selected = years.Y.selection.paper_rows / invalid_stratum_rows / leak.selected_events;
%   usable = years.Y.paper.events / invalid_stratum.events / leak.events (scripts/paper_propagation_tables.py).

\paragraph{RPKI validation.}
We validate the \emph{new} origin of every replayed origin change under RFC~6811~\cite{rfc6811}.
We use the ROAs that RIPE NCC archived for that UTC day~\cite{ripe_rpki_archive}, or the nearest earlier archive within one week.
The archive has no documented snapshot time, so a ROA change on the event day can give the wrong state.
% src (round 4, panel-r4 Q19, S5 "ROA snapshot time"): src/bgbench/rpki_validate.py roas_for (file of the event's UTC day, else the nearest
%   earlier day within 7 days); scripts/fetch_roas.py (one roas.csv.xz per TAL and day from https://ftp.ripe.net/rpki/<tal>.tal/YYYY/MM/DD/);
%   ftp.ripe.net/rpki/README.txt and https://github.com/RIPE-NCC/internet-dataset-descriptions/blob/main/rpki-repo-archive.md (both read
%   2026-10-31) state no snapshot time; docs/results/2026-09-09-rpki-validity.md caveat ("a ROA created hours after the event, or an event minutes
%   after a ROA change, can be misclassified"); docs/results/2026-09-15-decix-rs-dumps.md (disagreements with a live validator are archive-timing
%   artefacts of the daily snapshot).
% round 5 (panel-r5 NEW-10 / X11, domain B3): "all five trust anchors" is not guaranteed: scripts/fetch_roas.py (TALS ripencc, arin, apnic,
%   lacnic, afrinic) skips a trust anchor whose daily file returns HTTP 404 and writes the day file from the others; src/bgbench/rpki_validate.py
%   roas_for falls back to an earlier day only when the whole day file is missing. A missing anchor turns its invalid routes into unknown ones.
A revision analysis finds AFRINIC missing on 2026-04-04, and one anchor below 90\,\% of its local median size on 46 more days, none in R2.
% local median: the median of the same trust anchor over the 15 files on each side (scripts/paper_roa_anchor_check.py NEIGHBOURS, SHORT_RATIO).
These 47 days hold \num{1049} of \num{79870} events, and 12 of them change a label with the anchor of the nearest earlier or later day.
Both days agree on 3 changes from unknown to invalid, a fourth event changes differently with each, and 8 change with one day only.
% src (round 6, panel-r6 NEW-6 / X1; review-r6 W3, item 3): roa_anchor_check.json development.revalidation_on_flagged_files:
%   events_with_any_label_change 12 (sum of by_file[*].events_with_any_label_change); validity_changes_where_earlier_and_later_repairs_agree
%   {"unknown->invalid": 3}; assessment.new_origin_both_repairs_agree_and_differ {"unknown->invalid": 3}, displaced_origin_both_repairs_agree_and_differ
%   {"unknown->valid": 3}, stored_unknown_covered_under_both_repairs_but_class_differs 1 (20230804T010049_216.247.106.0-24_9009: earlier
%   [valid, invalid], later [invalid, valid]), one_repair_only_adjacent_day_roa_change 8 (3 + 1 + 8 = 12). The former "at most 3 labels" was not a
%   bound, and "in full" claimed more than a presence and row-count check shows: the printed text now names presence (tal_absent) and the
%   short-file criterion (rows < 0.9 x local median, "usual size") only. No claim of complete archives.
% src (round 5, panel-r5 A-R1(b), NEW-10 / X11): paper/generated/roa_anchor_check.json (scripts/paper_roa_anchor_check.py, first commit
%   b644cf6b 2026-10-31 17:47 UTC). development.files_with_tal_absent_from_first_used_file {"2026-04-04": ["afrinic"]}; ripe_archive_check
%   ripe_header_only ["2026-04-04 afrinic"] (empty in RIPE's own archive), stored_equals_ripe 47 of 47; event_days.days_resolving_to_flagged_file 47
%   (tal_absent 1, short 46: tal_0.5_to_0.9_of_local_median 40 + tal_below_0.5 6; "short" = rows < 0.9 x median of that TAL in the +-15 neighbouring
%   files; summary: "one anchor at 10%-85% of its usual size"); events_on_flagged_files.events 1,049; events 79,870 (reproduction: 79,870 identical).
%   "At most 3": revalidation_on_flagged_files.validity_changes_where_earlier_and_later_repairs_agree {"unknown->invalid": 3} (repair = the affected
%   TAL from the nearest unflagged earlier or later file within 7 days); summary_detail "Stored valid or invalid labels that change under both
%   repairs: none"; the 8 changes under one repair only are ROA churn on the adjacent day (assessment.one_repair_only_adjacent_day_roa_change).
%   "None in R2": newperiod.files_flagged {} and calendar_r2_window (45 days 2026-09-01..2026-10-15, all own unflagged files); all 47 new-period
%   files complete; 224,466 new-period labels reproduced exactly. The 6 fallback days (2023-06-24, 2023-07-14..17, 2026-06-13) are absent from
%   RIPE's archive for all five anchors (fallback_days_absent_at_ripe_for_all_tals), so the one-week fallback covers them.

Of the invalid events, 98--99\,\% are \texttt{invalid\_asn}, where no covering ROA names the new origin.
The others are \texttt{invalid\_length}, where the prefix exceeds the maxLength of a ROA for the origin.
We also validate the \emph{displaced} origin.
An invalid origin that displaces a valid one is an \emph{outward} change.
Its \emph{return} is the first identified event that restores the previous origin of the prefix within 24\,h (\Cref{sec:recovery}).
% src: paper/generated/propagation-facts.json (invalid_asn 8,354 / 8,492 in 2024 and 8,818 / 8,898 in 2025)

\paragraph{Replay and targets.}
We replay the window of each event and rebuild its anomalous announcements.
For an origin change, these carry the prefix with the new origin.
For a leak, their path runs through the new triple.
The onset $t_0$ is the first observation of an anomalous route at any collector.
The replay counts a route that a session reset or table transfer sends again like any other announcement.
% src (round 4, panel-r4 Q19 "session churn and table transfers"; review-r4 W5 "peer/session churn"): src/bgbench/replay_day.py and
%   src/bgbench/replay.py keep every announcement of the event prefix in the window and have no session-reset or table-dump handling;
%   src/bgbench/prefix_events.py compares each announcement with the previous origin of the same peer and prefix on the same collector-day.
%   The only reset hook (src/bgbench/forecast_signals.py reset_peer, "never infer it") belongs to the online trigger, which no reported result uses.

The \textbf{membership} set $\set{T}$ holds every AS on an anomalous path from a collector peer to the actor by $t_0+T$.
The actor is the new origin or the leaker.
For leaks we drop the hops beyond the leaker towards the origin.
The set accumulates: an AS stays in it once observed.

\paragraph{Persistence floor.}
Persistence forecasts that the set seen at horizon \Dl is the final set.
Its Jaccard with the two-hour set, $\Jac(\set{\Dl},\set{2h})$, is the \floor.
Because $\set{\Dl}\subseteq\set{2h}$, the floor equals $|\set{\Dl}|/|\set{2h}|$, the share of the two-hour set already visible at \Dl.
We use $\Dl=30$\,s unless we say otherwise.

This is close to the p90 collector delay of announcements, 26\,s (median 2\,s), which we measured with the RIS routing beacons~\cite{ripe_ris_routing_beacons} in June 2024.
Withdrawals reach the collectors later, with a median of 97\,s and a p90 of 270\,s.
It also equals the eBGP minimum route advertisement interval that RFC~4271 suggests~\cite{rfc4271}.
The 30-s floor can therefore partly reflect export timing at collector peers.
We report floors on the test split of the forecast experiments, September to December of each year (June to August in 2026).
% src: docs/paper/corrections-2026-10-16/beacon-delay/summary.json (announcement delay p50 2 s, p90 26 s)
